% Publication fragment for the existing BGS_unified entry point.
\subsection{The actual higher-seven-primary commutant}
\label{sec:actual-seven-primary-commutant}
For every odd $N$ satisfying $49\mid N$ and $5\nmid N$, the prescribed
pair $A_N,M_N$ has common commutant $\mathbb C I$ (FM-024).
This is an actual-model conclusion. Reducing $1000A_N$ at the Gaussian
prime with $i\mapsto2$ kills the diagonal and lower bonds, while every
internal upper bond and the raw metaplectic Fourier kernel remain units.
Both matrices are conjugated by the same chain-normalizing gauge.

The exact last and penultimate block-row equations force the leading
Toeplitz coefficient to be scalar. Their sufficient conditions are tied
by polynomial identities to six numerators and a degree-$828$ norm.
Complete factor products and Rabin tests exclude factors whose degrees
are divisible by $42=\operatorname{ord}_{49}(5)$, uniformly excluding
zeros at all roots of order divisible by $49$. Minimal-tail elimination
handles $N\ge147$; an exact actual-residue cyclic-rank certificate at
$N=49$ handles the short-chain base. A nonzero commutator-rank minor
then lifts to characteristic zero.

For an algebraic-entry matrix $X$, writing the algebraic Hermitian
spectral resolution $A_N=\sum_j a_jP_j$ in
$[X,e^{itA_N}M_N]=0$ gives
$\sum_j e^{ita_j}(XP_jM_N-P_jM_NX)=0$.
For nonzero algebraic real $t$, Lindemann--Weierstrass coefficient
separation~\cite{Popescu2023} yields each coefficient zero. Summing shows $[X,M_N]=0$,
then $[X,P_j]=0$, hence $[X,A_N]=0$. Thus $X$ is scalar, including at
$t=10^{-80}N$. This is only an algebraic-entry commuter restriction.

The higher-seven component excludes $5\mid N$ and $7\parallel N$; the separate first-seven extension follows below. No spectral
simplicity, quantitative gap, signed Heisenberg comparison, or BGS
conclusion is asserted. The canonical proof, completed independent
review, and portable exact certificate replays are separate from the
bounded Lean predicate and its native-evaluation trust boundary.

\paragraph{Finite formal certificate.} The seven-primary Lean source has compiled proofs of 22 executable array predicates, with independent semantic review and coefficient binding. These use \texttt{native\_decide}; generic polynomial correctness, Rabin soundness, finite-field bridges and the commutant proof remain unformalized.

\paragraph{Extension to the first unramified seven-primary level.}
A separate independently reviewed exact certificate closes the remaining
unramified case: for odd $N>7$ with $7\mid N$, $49\nmid N$ and $5\nmid N$,
the actual pair $(A_N,M_N)$ has scalar complex joint commutant.
Four corner-block equations eliminate the long endpoint gauge and give a
$63\times48$ polynomial system over $\mathbb F_{5^6}$. In each of six
residue classes, three minors have degree at most $936$. Their $937$-node
identities, plus ten extra nodes each, are independently rebuilt in
$17{,}046$ exact determinant evaluations. Gcd, Bezout and finite-field
odd-order exclusions then establish the rank for every relevant root;
this is not dimensional sampling. Invertibility of the corner block removes
all Toeplitz tails, and a nonzero commutator minor lifts to characteristic zero.
Together with the higher-seven result this extends FM-024 to all odd $N>7$
with $7\mid N$ and $5\nmid N$. Combining with the separate prior FM-008
result covers odd $N\ge3$ with $35\nmid N$, except $N=7$.
At that stage the remaining ramified family was the odd multiples of $35$; the chain-trace extension below closes its first-seven part.
This does not make the full complex commutant of $U_N$ scalar, prove
seven-divisible spectral simplicity or establish BGS. FORMAL-003 remains
restricted to its earlier executable-array predicates.

\paragraph{First-seven closure with arbitrary five-adic ramification.}
A separate independently reviewed argument removes the condition $5\nmid N$
at the first seven-primary level: the actual pair $(A_N,M_N)$ has scalar
complex joint commutant for every odd $N>7$ with $7\mid N$ and $49\nmid N$.
A valuation gauge preserves a rank-seven endpoint block even when the full
reduced Fourier matrix is singular. Before reduction, an integral block-trace
consequence of the actual commutator supplies an additional diagonal constraint.
Its two normalized harmonics cannot vanish together, by an exact polynomial
Bezout identity. All nine collision partitions have enough singleton groups;
nonzero Fourier minors of orders at most three force every block scalar.
A nonzero minor of the augmented integral system then lifts to characteristic
zero. The finite controls support this uniform argument, not dimensional
extrapolation. Together with the separate higher-seven and FM-008 results,
the then-remaining uniform arithmetic family was the odd multiples of $245$,
closed by the separate corrected lattice argument below,
besides the genuine exception $N=7$. No single-$U$ complex scalar-commutant,
seven-divisible simplicity, quantitative gap or BGS conclusion follows.
